%% =============================================================================
%% BHS Counterexample Paper — Minimum Viable LaTeX Template
%% Target: arXiv (math.CO) + Linear Algebra and its Applications (Elsevier)
%% Document class: elsarticle (Elsevier standard)
%%
%% HOW TO USE:
%%   Overleaf: Upload this file + references.bib -> Recompile
%%   Local:    pdflatex main && bibtex main && pdflatex main && pdflatex main
%% =============================================================================

\documentclass[preprint,12pt]{elsarticle}

%% --- Essential Packages ---
\usepackage{amsmath}      % align, cases, equation*, etc.
\usepackage{amssymb}      % \mathbb, \leq, etc.
\let\proof\relax \let\endproof\relax
\usepackage{amsthm}       % theorem, proof environments
\usepackage{booktabs}     % Professional tables: \toprule, \midrule, \bottomrule
\usepackage{graphicx}     % \includegraphics for figures
\usepackage{hyperref}     % Clickable cross-references and URLs
\usepackage{enumitem}     % Better list formatting
\usepackage{mathtools}    % \coloneqq and other math extras

%% --- Theorem-like Environments ---
%% Numbers are shared and sequential within each section: Theorem 2.1, Lemma 2.2, ...
\newtheorem{theorem}{Theorem}[section]
\newtheorem{lemma}[theorem]{Lemma}
\newtheorem{proposition}[theorem]{Proposition}
\newtheorem{corollary}[theorem]{Corollary}
\newtheorem{conjecture}[theorem]{Conjecture}

\theoremstyle{definition}
\newtheorem{definition}[theorem]{Definition}
\newtheorem{example}[theorem]{Example}

\theoremstyle{remark}
\newtheorem{remark}[theorem]{Remark}

%% --- Custom Commands (save typing, ensure consistency) ---
\newcommand{\SoC}{\mathrm{SoC}}       % StarOfCliques
\newcommand{\SoS}{\mathrm{SoS}}       % StarOfStars
\newcommand{\DS}{\mathrm{DS}}          % DoubleStar
\newcommand{\Cat}{\mathrm{Cat}}        % Caterpillar
\newcommand{\Lspec}{\mu}               % Laplacian spectral radius

%% --- Hyperref Setup ---
\hypersetup{
  colorlinks=true,
  linkcolor=blue,
  citecolor=blue,
  urlcolor=blue,
}

%% --- Remove elsarticle default footer (optional, cleaner for arXiv) ---
\journal{Linear Algebra and its Applications}

%% =============================================================================
%% DOCUMENT BEGINS
%% =============================================================================
\begin{document}

\begin{frontmatter}

%% --- Title ---
\title{Counterexamples to Five Conjectured Upper Bounds\\
  for the Largest Laplacian Eigenvalue}

%% --- Author ---
\author[inst1]{Taewoo Ha}
\ead{txh2120@gmail.com}

\affiliation[inst1]{
  organization={Independent Researcher},
  city={California},
  country={United States}
}

%% --- Abstract ---
\begin{abstract}
  Brankov, Hansen, and Stevanovi\'{c} (2006) generated 68 conjectured
  upper bounds for the largest Laplacian
  eigenvalue~$\mu(G)$ of a graph~$G$. Ghebleh et al.\ (2026) disproved
  30 of these, leaving 38 open. We resolve five of the remaining open
  conjectures (Bounds~11, 13, 40, 45, and~48) by constructing explicit
  counterexample families---star-of-cliques and double star
  graphs---whose Laplacian spectral radii we compute analytically via
  equitable partition quotient matrices. We identify hub--periphery degree
  imbalance as the structural mechanism underlying all violations and
  classify the remaining 33~open bounds by their resistance to this mechanism.
\end{abstract}

%% --- Keywords and MSC ---
\begin{keyword}
  Laplacian eigenvalue \sep spectral radius \sep upper bound \sep
  counterexample \sep equitable partition
  \MSC[2020] 05C50 \sep 15A18
\end{keyword}

\end{frontmatter}

%% =============================================================================
%% 1. INTRODUCTION
%% =============================================================================
\section{Introduction}\label{sec:intro}

Let $G = (V, E)$ be a simple connected graph on $n = |V|$ vertices and
$m = |E|$ edges. The \emph{Laplacian matrix} of~$G$ is
$L(G) = D(G) - A(G)$, where $A(G)$ is the adjacency matrix and
$D(G) = \operatorname{diag}(d_1, \ldots, d_n)$ is the diagonal matrix of
vertex degrees. Since $L(G)$ is positive semidefinite with row sums zero,
its eigenvalues satisfy $0 = \lambda_n \leq \cdots \leq \lambda_1$, and
the largest eigenvalue
\begin{equation}\label{eq:mu-def}
  \Lspec(G) \coloneqq \lambda_{\max}\bigl(L(G)\bigr)
\end{equation}
is called the \emph{Laplacian spectral radius}. Upper bounds on
$\Lspec(G)$ in terms of easily computed graph invariants are of both
theoretical and algorithmic interest; see~\cite{BrouwerHaemers2012} for
background. In 2006, Brankov, Hansen, and
Stevanovi\'{c}~\cite{BHS2006} generated 68~conjectured upper bounds
for~$\Lspec(G)$ using a backward-reconstruction algorithm that builds
expressions evaluating to~$2x$ and substitutes~$d_v$ or~$m_v$
(see~\cite{BHS2006} for details). These bounds are expressed solely
in terms of vertex degrees~$d_v$ and
\emph{average neighbor degrees}
\begin{equation}\label{eq:mv-def}
  m_v = \frac{1}{d_v} \sum_{u \sim v} d_u.
\end{equation}
The 68~conjectures fall into two families.
\emph{Vertex-type} bounds have the form
$\Lspec(G) \leq \max_v f(d_v, m_v)$, and \emph{edge-type}
bounds have the form
\[
  \Lspec(G) \leq \max_{ij\in E} g(d_i, d_j, m_i, m_j).
\]
None had been proved or disproved for nearly two decades.

Ghebleh, Al-Yakoob, Kanso, and Stevanovi\'{c}~\cite{Ghebleh2024}
recently applied the cross-entropy method (CEM) of
Wagner~\cite{Wagner2021} to disprove 30 of the 68~bounds by
constructing small counterexample graphs, leaving 38~bounds open. Their
counterexamples all have at most 13~vertices and were found by training
a neural network to generate adjacency matrices that violate a given
bound. Subsequently, Bouffard and Breen~\cite{BouffardBreen2025}
parallelized the CEM approach and extended the search to larger graphs,
but reported no new counterexamples among the remaining 38~bounds.

In the present paper we disprove five of the 38~remaining open bounds.

\begin{theorem}[Main result]\label{thm:main}
  BHS Bounds~11, 13, 40, 45, and~48 from~\cite{BHS2006} are false.
  Specifically:
  \begin{enumerate}[label=\textup{(\alph*)}]
    \item Bound~11 is violated by $\SoC(K_m, t)$ for all $m \geq 4$
      at $t = 2m - 1$
      \textup{(Theorem~\ref{thm:bound11})}.
    \item Bound~13 is violated by $\SoC(K_m, t)$ for all $m \geq 5$
      at $t = 2m - 1$
      \textup{(Theorem~\ref{thm:bound13})}.
    \item Bound~40 is violated by the star of stars
      $\SoS(39, 8)$ and $\SoS(49, 9)$
      \textup{(Theorem~\ref{thm:bound40})}.
    \item Bound~45 is violated by $\DS(k,k)$ for all $k \geq 6$
      \textup{(Theorem~\ref{thm:bound45})}.
    \item Bound~48 is violated by the path~$P_3$
      \textup{(Remark~\ref{rem:bound48})}.
  \end{enumerate}
\end{theorem}

\noindent
In each case the proof is constructive: we exhibit an infinite family
(or, for Bound~48, a single small graph) together with an analytic
computation of its Laplacian spectral radius via equitable partition
quotient matrices.

Our approach differs fundamentally from the optimization-based methods
of~\cite{Ghebleh2024,Wagner2021,BouffardBreen2025}. Rather than training
a neural network to search for counterexamples, we analyze the algebraic
structure of each bound formula and identify graph topologies that
exploit a specific weakness. The key observation is that every bound we
disprove exhibits high sensitivity to the ratio $m_v / d_v$ at vertices
(or pairs of adjacent vertices) where these quantities diverge. Graphs
with \emph{hub--periphery} structure---a high-degree hub connected to
dense peripheral subgraphs---create exactly this divergence: the hub's
average neighbor degree~$m_v$ is determined by the peripheral structure
and can be made large relative to the hub degree~$d_v$, or vice versa.
Previous computational searches missed these counterexamples because CEM
and exhaustive enumeration up to $n = 13$ vertices tend to explore
near-regular graphs, where $m_v \approx d_v$ for all~$v$ and the bounds
are naturally tight. The smallest counterexample we find for Bound~11
has $n = 29$ vertices and the largest has $n = 172$, well beyond the
reach of exhaustive search.

The paper is organized as follows. Section~\ref{sec:prelim} fixes
notation and recalls the theory of equitable partitions. It also states
the five BHS bounds that we disprove. Section~\ref{sec:families} defines
the counterexample graph families: star of cliques, double star, and star
of stars. Section~\ref{sec:main} contains the main results: closed-form
spectral radius computations and the violation proofs for each of the
five bounds. Section~\ref{sec:mechanism} analyzes the structural
mechanisms underlying all five violations.
Section~\ref{sec:classification} classifies the remaining 33~open bounds
by their resistance to this mechanism. Section~\ref{sec:discussion}
discusses implications and directions
for future work.


%% =============================================================================
%% 2. PRELIMINARIES
%% =============================================================================
\section{Preliminaries}\label{sec:prelim}

\subsection{Graph notation}\label{ssec:notation}

Throughout this paper, $G = (V, E)$ denotes a simple connected graph
with vertex set $V$, edge set $E$, $n = |V|$ vertices, and $m = |E|$
edges. The \emph{adjacency matrix} $A(G)$ is the $n \times n$ matrix
with $A_{ij} = 1$ if $\{i,j\} \in E$ and $A_{ij} = 0$ otherwise. The
\emph{degree matrix} is $D(G) = \operatorname{diag}(d_1, \ldots, d_n)$,
where $d_v = \sum_{u} A_{vu}$ is the degree of vertex~$v$. The
\emph{Laplacian matrix} $L(G) = D(G) - A(G)$ is positive semidefinite
with smallest eigenvalue~$0$. We write
$\Lspec(G) = \lambda_{\max}(L(G))$ for the \emph{Laplacian spectral
radius}.

For each vertex $v \in V$ with $d_v \geq 1$, the \emph{average neighbor
degree} is
\[
  m_v = \frac{1}{d_v} \sum_{u \sim v} d_u,
\]
where the sum is over all neighbors~$u$ of~$v$. Note that $m_v \geq 1$
for every vertex in a connected graph, and $m_v = d_v$ for every vertex
if and only if $G$ is regular.

\subsection{Equitable partitions}\label{ssec:equitable}

Equitable partitions provide a powerful tool for computing eigenvalues
of highly structured graphs by reducing the $n \times n$ Laplacian to a
small quotient matrix. We recall the essential definitions; for a
comprehensive treatment see Godsil and Royle~\cite{GodsilRoyle2001}.

\begin{definition}[Equitable partition]\label{def:equitable}
  A partition $\pi = \{C_1, \ldots, C_k\}$ of the vertex set~$V(G)$ is
  \emph{equitable} if, for every pair $(i,j)$ with
  $1 \leq i, j \leq k$, every vertex in~$C_i$ has the same number of
  neighbors in~$C_j$. We denote this common value by~$b_{ij}$.
\end{definition}

Given an equitable partition with cells $C_1, \ldots, C_k$, the
\emph{quotient matrix} is the $k \times k$ matrix $Q$ with entries
$Q_{ij} = b_{ij}$, and the \emph{quotient Laplacian} is the $k \times k$
matrix $Q_L$ with entries
\[
  (Q_L)_{ij} =
  \begin{cases}
    d_{C_i}  & \text{if } i = j, \\
    -b_{ij}  & \text{if } i \neq j,
  \end{cases}
\]
where $d_{C_i} = \sum_{j=1}^{k} b_{ij}$ is the common degree of every
vertex in cell~$C_i$.

\begin{theorem}[{\cite[Theorem~9.3.3]{GodsilRoyle2001}}]\label{thm:equitable}
  Let $\pi$ be an equitable partition of~$G$ with quotient Laplacian
  $Q_L$. Then every eigenvalue of $Q_L$ is an eigenvalue of~$L(G)$.
  In particular,
  \[
    \Lspec(G) \geq \lambda_{\max}(Q_L),
  \]
  with equality whenever the eigenvector of~$L(G)$ corresponding to
  $\Lspec(G)$ is constant on each cell of~$\pi$.
\end{theorem}

\noindent
In all counterexample families considered in this paper, the natural
partition by vertex type is equitable, and the largest eigenvalue
eigenvector is piecewise constant on cells. This reduces the computation
of~$\Lspec(G)$ to finding the largest root of a polynomial of degree at
most~$3$.

\subsection{The BHS bounds}\label{ssec:bhs-bounds}

We now state the five conjectured bounds from~\cite{BHS2006} that we
disprove. All five are listed as open in~\cite[Table~1]{Ghebleh2024}.
Bounds~11 and~13 are vertex-type bounds; Bounds~40, 45, and~48 are
edge-type bounds.

\begin{conjecture}[BHS Bound~11 {\cite{BHS2006}}]\label{conj:bound11}
  For every connected graph~$G$,
  \[
    \Lspec(G) \leq \max_{v \in V} \frac{2\, m_v^{\,3}}{d_v^{\,2}}.
  \]
\end{conjecture}

\begin{conjecture}[BHS Bound~13 {\cite{BHS2006}}]\label{conj:bound13}
  For every connected graph~$G$,
  \[
    \Lspec(G) \leq \max_{v \in V} \frac{2\, m_v^{\,4}}{d_v^{\,3}}.
  \]
\end{conjecture}

\begin{conjecture}[BHS Bound~40 {\cite{BHS2006}}]\label{conj:bound40}
  For every connected graph~$G$,
  \[
    \Lspec(G) \leq \max_{(i,j) \in E}\; \Bigl(2 + \sqrt{2(m_i{-}1)^2 + 2(m_j{-}1)^2
      + d_i^2 + d_j^2 - d_i m_i - d_j m_j}\,\Bigr).
  \]
\end{conjecture}

\begin{conjecture}[BHS Bound~45 {\cite{BHS2006}}]\label{conj:bound45}
  For every connected graph~$G$,
  \[
    \Lspec(G) \leq \max_{(i,j) \in E}\; \Bigl(2 + \sqrt{(d_i{-}d_j)^2 + 2(d_i m_i + d_j m_j)
      - 4(m_i + m_j) + 4}\,\Bigr).
  \]
\end{conjecture}

\begin{conjecture}[BHS Bound~48 {\cite{BHS2006}}]\label{conj:bound48}
  For every connected graph~$G$,
  \[
    \Lspec(G) \leq \max_{(i,j) \in E}\;
      \frac{2(d_i^2 + d_j^2)}{2 + \sqrt{2(d_i^2 + d_j^2) - 4(m_i + m_j) + 4}}.
  \]
\end{conjecture}

\noindent
Bounds~11 and~13 are structurally similar: both take the form
$\max_v 2m_v^p/d_v^{p-1}$ with $p = 3$ and $p = 4$,
respectively.  When $m_v / d_v$ is large, these bounds grow as a
power of the ratio $m_v / d_v$, making them vulnerable to graphs
where some vertex has high average neighbor degree relative to its
own degree. Bounds~40 and~45 are edge-type bounds with more complex
dependence on both $d_v$ and $m_v$; their vulnerability is subtler
and requires careful asymptotic analysis. Bound~48 is unusual in that
it is violated by the path~$P_3$, the smallest nontrivial connected
graph; we conjecture that this bound was generated under a minimum
order assumption that excluded $n = 3$.


%% =============================================================================
%% 3. COUNTEREXAMPLE GRAPH FAMILIES
%% =============================================================================
\section{Counterexample graph families}\label{sec:families}

\begin{definition}[Star-of-Cliques]\label{def:soc}
  For integers $m \geq 3$ and $t \geq 2$, the graph
  $\SoC(K_m, t)$ consists of a central \emph{hub} vertex~$v_0$
  connected by a single edge to one designated vertex (the
  \emph{connector}) in each of $t$ disjoint copies of the complete
  graph~$K_m$. The total number of vertices is $n = 1 + tm$.
\end{definition}

\begin{definition}[Double Star]\label{def:ds}
  For an integer $k \geq 1$, the double star $\DS(k,k)$ is obtained by
  connecting two vertices $u$ and~$v$, each of which is adjacent
  to~$k$ pendant vertices. The total number of vertices is $n = 2k + 2$.
\end{definition}

\begin{definition}[Star-of-Stars]\label{def:sos}
  For integers $t \geq 2$ and $s \geq 1$, the graph $\SoS(t, s)$
  consists of a central \emph{hub} vertex~$v_0$ adjacent to $t$
  \emph{intermediate} vertices $w_1, \ldots, w_t$, each of which is
  adjacent to $s$ pendant \emph{leaf} vertices. The total number of
  vertices is $n = 1 + t + ts = 1 + t(1 + s)$.
\end{definition}

\begin{definition}[Caterpillar]\label{def:cat}
  For integers $\ell \geq 2$ and $k \geq 0$, the caterpillar
  $\Cat(\ell, k)$ is obtained from a path on $\ell$ vertices (the
  \emph{spine}) by attaching $k$ pendant leaves to each spine vertex.
  The total number of vertices is $n = \ell(k + 1)$.
\end{definition}


%% =============================================================================
%% 4. MAIN RESULTS
%% =============================================================================
\section{Main results}\label{sec:main}

%% --- Example: Spectral radius of StarOfCliques ---

\begin{proposition}\label{prop:mu-soc}
  For $m \geq 3$ and $t \geq 2$, the Laplacian spectral radius of
  $\SoC(K_m, t)$ is
  \begin{equation}\label{eq:mu-soc}
    \Lspec\bigl(\SoC(K_m, t)\bigr) =
    \frac{(m+t+1) + \sqrt{(t-m)^2 + 2(t+m) - 3}}{2}.
  \end{equation}
\end{proposition}

\begin{proof}
  The vertex set admits an equitable partition into three cells:
  $H = \{v_0\}$ (hub), $C$ (connectors, $|C| = t$), and $I$ (interior,
  $|I| = t(m-1)$). The quotient Laplacian is the $3 \times 3$ matrix
  \begin{equation}\label{eq:quotient-soc}
    Q_L =
    \begin{pmatrix}
      t    & -t     & 0      \\
      -1   & m      & -(m-1) \\
      0    & -1     & 1
    \end{pmatrix}.
  \end{equation}
  One eigenvalue is $\lambda_1 = 0$ (since the quotient Laplacian has zero row sums). The
  remaining two eigenvalues are roots of
  \[
    \lambda^2 - (m + t + 1)\lambda + (tm + 1) = 0,
  \]
  giving the largest eigenvalue as~\eqref{eq:mu-soc}. Since the
  corresponding eigenvector is piecewise constant on the cells,
  this is also the largest eigenvalue of~$L(G)$.
\end{proof}


%% --- Example: Bound 11 violation theorem ---

\begin{theorem}\label{thm:bound11}
  BHS Bound~11,
  \begin{equation}\label{eq:bound11}
    \Lspec(G) \leq \max_{v \in V} \frac{2m_v^3}{d_v^2},
  \end{equation}
  is violated by $\SoC(K_m, t)$ for all $m \geq 4$ at $t = 2m - 1$.
  In particular, $\SoC(K_8, 15)$ with $n = 121$ satisfies
  $\Lspec(G) \approx 16.796$ while the bound gives $\approx 16.000$.
\end{theorem}

\begin{proof}
  By Proposition~\ref{prop:mu-soc}, we have
  $\Lspec = \frac{(m+t+1) + \sqrt{(t-m)^2 + 2(t+m)-3}}{2}$.
  The hub vertex~$v_0$ has $d_{v_0} = t$ and $m_{v_0} = m$ (all
  neighbors are connectors of degree~$m$), so Bound~11 evaluates to
  \[
    \frac{2m_{v_0}^3}{d_{v_0}^2} = \frac{2m^3}{t^2}.
  \]
  For large~$t$, this is $O(t^{-2})$ while $\Lspec$ grows as
  $t + O(1)$, so the violation eventually occurs. Specifically,
  at $t = 2m - 1$ the connector vertex has $d = m$ and
  $m_v = \bigl(t + (m-1)^2\bigr)/m = m$, giving
  $\frac{2m_v^3}{d^2} = 2m$, which dominates the hub bound.
  By Proposition~\ref{prop:mu-soc},
  $\Lspec = \bigl(3m + \sqrt{m^2 + 4m - 4}\bigr)/2 > 2m$
  since $\sqrt{m^2 + 4m - 4} > m$ for $m \geq 2$ (obtained by
  squaring). The violation persists in a neighborhood of $t = 2m - 1$ whose width
  grows with~$m$ (computations show windows $[5,7]$ for $m=4$ up to
  $[22,137]$ for $m=12$), but for fixed~$m$ the connector bound
  $2m_v^3/m^2$ eventually grows as $\Theta(t^3/m^2)$, exceeding
  the linearly growing $\Lspec$, so the violation window is finite.
\end{proof}


%% --- Bound 13 violation ---

\begin{theorem}\label{thm:bound13}
  BHS Bound~13,
  \begin{equation}\label{eq:bound13}
    \Lspec(G) \leq \max_{v \in V} \frac{2m_v^4}{d_v^3},
  \end{equation}
  is violated by $\SoC(K_m, t)$ for all $m \geq 5$ and $t = 2m - 1$.
  In particular, $\SoC(K_9, 17)$ with $n = 154$ satisfies
  $\Lspec(G) \approx 18.815$ while the bound gives $\approx 18.000$.
\end{theorem}

\begin{proof}
  The connector vertex has $d = m$ and
  $m_v = \bigl(t + (m-1)^2\bigr)/m$. At $t = 2m - 1$ we have
  $t + (m-1)^2 = m^2$, giving
  $B_{13}^{\mathrm{conn}} = 2m^4/m^3 = 2m$.
  By Proposition~\ref{prop:mu-soc}, $\Lspec = (3m + \sqrt{m^2 + 4m - 4})/2$.
  Since $\sqrt{m^2 + 4m - 4} > m$ for $m \geq 2$, we obtain
  $\Lspec > 2m$, establishing the violation for $m \geq 6$.
  For $m = 5$, the interior bound is $\approx 10.195$ but
  $\Lspec\bigl(\SoC(K_5, 9)\bigr) \approx 10.702 > 10.195$.
\end{proof}


%% --- Spectral radius of DoubleStar ---

\begin{proposition}\label{prop:mu-ds}
  For $k \geq 2$, the Laplacian spectral radius of $\DS(k,k)$ is
  \begin{equation}\label{eq:mu-ds}
    \Lspec\bigl(\DS(k,k)\bigr)
    = \frac{(k+3) + \sqrt{k^2 + 6k + 1}}{2}.
  \end{equation}
\end{proposition}

\begin{proof}
  The $\mathbb{Z}_2$ symmetry swapping the two hubs decomposes
  eigenspaces into symmetric and antisymmetric sectors.

  \emph{Symmetric sector.} Quotient matrix
  $Q_+ = \bigl(\begin{smallmatrix} k & -k \\ -1 & 1 \end{smallmatrix}\bigr)$,
  eigenvalues $0$ and $k+1$.

  \emph{Antisymmetric sector.} Quotient matrix
  $Q_- = \bigl(\begin{smallmatrix} k+2 & -k \\ -1 & 1 \end{smallmatrix}\bigr)$,
  characteristic polynomial $\lambda^2 - (k+3)\lambda + 2 = 0$,
  eigenvalues $(k+3 \pm \sqrt{k^2+6k+1})/2$.

  Pendant leaves contribute eigenvalue~$1$ with multiplicity $2(k-1)$.
  Since $\sqrt{k^2+6k+1} > k-1$ for $k \geq 2$, the antisymmetric
  $\lambda_+$ is the global maximum.
\end{proof}


%% --- Bound 45 violation ---

\begin{theorem}\label{thm:bound45}
  BHS Bound~45,
  \begin{equation}\label{eq:bound45}
    \Lspec(G) \leq \max_{(i,j) \in E}\; \Bigl(2 + \sqrt{(d_i{-}d_j)^2 + 2(d_i m_i + d_j m_j)
      - 4(m_i{+}m_j) + 4}\,\Bigr),
  \end{equation}
  is violated by $\DS(k,k)$ for all $k \geq 6$.
  At $k = 6$, $\DS(6,6)$ with $n = 14$ satisfies
  $\Lspec(G) \approx 8.772$ while the bound gives $\approx 8.676$.
\end{theorem}

\begin{proof}
  The hub--leaf edge maximizes Bound~45 for $k \geq 6$. We show
  $\Lspec > B_{45}^{\mathrm{hl}}$. Both $\Lspec - 2$ and
  $B_{45}^{\mathrm{hl}} - 2$ are positive for $k \geq 2$, so
  $\Lspec > B_{45}^{\mathrm{hl}}$ if and only if
  $(\Lspec - 2)^2 > (B_{45}^{\mathrm{hl}} - 2)^2$.
  Expanding and simplifying, then squaring a second time to
  clear the remaining radical, reduces to
  $4k(k^3 + 8k^2 - 11k - 2) > 0$, which holds for $k \geq 2$ since
  $k^3 + 8k^2 - 11k - 2 > 0$ for $k \geq 2$.
  The gap satisfies $B_{45} - \Lspec = -1/(2k) + O(k^{-2})$.
\end{proof}


%% --- Spectral radius of StarOfStars ---

\begin{proposition}\label{prop:mu-sos}
  For $t \geq 2$ and $s \geq 1$, the Laplacian spectral radius of
  $\SoS(t, s)$ is
  \begin{equation}\label{eq:mu-sos}
    \Lspec\bigl(\SoS(t, s)\bigr)
    = \frac{(t + s + 2) + \sqrt{(t - s)^2 + 4s}}{2}.
  \end{equation}
\end{proposition}

\begin{proof}
  Equitable partition into $\{$hub$\}$, $\{$intermediates$\}$,
  $\{$leaves$\}$. The quotient Laplacian is
  $Q_L = \bigl(\begin{smallmatrix} t & -t & 0 \\ -1 & 1{+}s & -s \\ 0 & -1 & 1 \end{smallmatrix}\bigr)$
  with characteristic polynomial
  $-\lambda[\lambda^2 - (t+s+2)\lambda + n]$, where $n = 1 + t(1+s)$.
  Discriminant $(t-s)^2 + 4s > 0$.
\end{proof}


%% --- Bound 40 violation ---

\begin{theorem}\label{thm:bound40}
  BHS Bound~40,
  \begin{equation}\label{eq:bound40}
    \Lspec(G) \leq \max_{(i,j) \in E}\; \Bigl(2 + \sqrt{2(m_i{-}1)^2 + 2(m_j{-}1)^2
      + d_i^2 + d_j^2 - d_i m_i - d_j m_j}\,\Bigr),
  \end{equation}
  is violated by $\SoS(39, 8)$ with $n = 352$, where
  $\Lspec(G) \approx 40.256$ and the bound gives $\approx 38.982$.
\end{theorem}

\begin{proof}
  The hub--intermediate edge has $d_h = 39$, $m_h = 9$,
  $d_w = 9$, $m_w = 47/9$. By Proposition~\ref{prop:mu-sos},
  $\Lspec = (49 + \sqrt{993})/2$. The inequality
  $\Lspec > B_{40}$ reduces to $\sqrt{993} > 27.25$, i.e.,
  $993 > 742.6$, which holds.
  Similarly, $\SoS(49, 9)$ with $n = 491$ gives
  $\Lspec = (60 + \sqrt{1636})/2 \approx 50.224$
  while $B_{40} \approx 48.487$ at the hub--intermediate edge
  (with $d_h = 49$, $m_h = 10$, $d_w = 10$, $m_w = 58/10$),
  confirming the violation with gap $\approx 1.74$.
\end{proof}


%% --- Bound 48 violation ---

\begin{remark}[Bound~48 --- degenerate case]\label{rem:bound48}
  BHS Bound~48,
  \[
    \Lspec(G) \leq \max_{(i,j) \in E}\;
    \frac{2(d_i^2 + d_j^2)}{2 + \sqrt{2(d_i^2 + d_j^2) - 4(m_i + m_j) + 4}},
  \]
  is violated by~$P_3$: the bound gives $10/(2{+}\sqrt{2}) \approx 2.929$
  while $\Lspec(P_3) = 3$. This is the only counterexample found.
\end{remark}


%% --- Summary Table ---

\begin{table}[t]
  \centering
  \caption{Summary of counterexamples to BHS bounds. ``Gap'' denotes
    the bound value minus~$\Lspec(G)$ (negative = violation).}
  \label{tab:counterexamples}
  \begin{tabular}{clcrrr}
    \toprule
    Bound & Counterexample & $n$ & $\Lspec(G)$ & Bound value & Gap \\
    \midrule
    11 & $\SoC(K_8, 15)$     & 121 & 16.7958 & 16.0000 & $-0.7958$ \\
    13 & $\SoC(K_9, 17)$     & 154 & 18.8151 & 18.0000 & $-0.8151$ \\
    40 & $\SoS(39, 8)$       & 352 & 40.2560 & 38.9818 & $-1.2741$ \\
    45 & $\DS(6, 6)$         &  14 &  8.7720 &  8.6762 & $-0.0958$ \\
    48 & $P_3$               &   3 &  3.0000 &  2.9289 & $-0.0711$ \\
    \bottomrule
  \end{tabular}
\end{table}



%% =============================================================================
%% 5. MECHANISM ANALYSIS
%% =============================================================================
\section{Why these bounds fail}\label{sec:mechanism}

The five bounds disproved in Section~\ref{sec:main} fail for related but
distinct structural reasons.  In this section we identify two mechanisms---hub
invisibility (vertex-type) and degree--spectrum decoupling (edge-type)---and
show that together they account for all five violations.  We then classify all
24~vertex-type BHS bounds by their hub behavior, providing a unified framework
that explains why the remaining vertex-type bounds resist our counterexample
families.

\subsection{Hub invisibility in vertex-type bounds}\label{ssec:hub-invisible}

Every vertex-type BHS bound has the form
$\Lspec(G) \leq \max_{v \in V} f(d_v, m_v)$ for some function~$f$.
On a star-of-cliques $\SoC(K_m, t)$ the hub vertex has
$(d_v, m_v) = (t, m)$, the connectors have
$(d_v, m_v) = \bigl(m, \, (t + (m{-}1)^2)/m\bigr)$,
and the interior vertices have
$(d_v, m_v) = \bigl(m{-}1, \, (m{-}2) + m/(m{-}1)\bigr)$.
The Laplacian spectral radius satisfies $\Lspec \sim t + 1$ as
$t \to \infty$ (Proposition~\ref{prop:mu-soc}), growing linearly in~$t$.

\begin{definition}[Hub-invisible bound]\label{def:hub-invisible}
  A vertex-type bound $\max_v f(d_v,m_v)$ is
  \emph{hub-invisible} on $\SoC(K_m, t)$
  if $f(t,m) \to 0$ as $t \to \infty$ with $m$ fixed,
  and \emph{hub-visible} if $f(t,m) \to \infty$.
\end{definition}

The following proposition classifies all~24 vertex-type BHS bounds.

\begin{proposition}[Classification of vertex-type bounds]\label{prop:hub-classification}
  Among the 24~vertex-type BHS bounds (Bounds~1, 4--14, 16, 18--27, 30),
  exactly two are hub-invisible:
  \begin{enumerate}[label=\textup{(\alph*)}]
    \item Bound~11, $f_{11}(d,m) = 2m^3/d^2$: hub contribution
      $f_{11}(t,m) = 2m^3/t^2 \to 0$.
    \item Bound~13, $f_{13}(d,m) = 2m^4/d^3$: hub contribution
      $f_{13}(t,m) = 2m^4/t^3 \to 0$.
  \end{enumerate}
  The remaining 22~vertex-type bounds are all hub-visible: for each,
  $f(t,m) \to \infty$ as $t \to \infty$.
\end{proposition}

\begin{proof}
  We verify the hub contribution $f(t,m)$ for each bound as $t \to \infty$
  with $m$ fixed.  Representative cases:
  \begin{itemize}
    \item \emph{Pure $d$-dominant bounds} (1, 4, 7, 8, 9, 14, 16):
      these satisfy $f(t,m) \geq c \cdot t^{\alpha}$ for some
      $\alpha \geq 1$ and $c > 0$.
      For example, $f_4(t,m) = 2t^2/m \to \infty$
      and $f_{14}(t,m) = 2t^3/m^2 \to \infty$.
    \item \emph{Mixed bounds} (5, 6, 10, 12, 18--27, 30): each contains
      at least one term growing in~$t$.
      For example, $f_5(t,m) = t^2/m + m \to \infty$;
      $f_6(t,m) = \sqrt{m^2 + 3t^2} \to \infty$;
      $f_{30}(t,m) = m^3/t^2 + t^2/m \to \infty$
      since the second term dominates.
    \item \emph{Hub-invisible bounds} (11, 13): these are the only bounds
      where every term in~$f$ has $d$ in the denominator with higher power
      than in the numerator.
      $f_{11}(t,m) = 2m^3/t^2 \to 0$;
      $f_{13}(t,m) = 2m^4/t^3 \to 0$.
  \end{itemize}
  The full evaluation for all~24 bounds is recorded in
  Table~\ref{tab:hub-classification}.
\end{proof}

\begin{table}[t]
  \centering
  \caption{Hub behavior of all 24 vertex-type BHS bounds on $\SoC(K_m, t)$
    as $t \to \infty$ with $m$ fixed.  Hub contribution $f(t,m)$ and its
    limit determine whether the bound is hub-visible~(V) or
    hub-invisible~(I).  Status:
    \textsc{Disproved}~=~disproved in Section~\ref{sec:main};
    \textsc{Open}~=~unresolved.}
  \label{tab:hub-classification}
  \begin{tabular}{clccc}
    \toprule
    Bound & Formula $f(d,m)$ & $f(t,m)$ & Hub & Status \\
    \midrule
    1  & $\sqrt{4d^3/m}$           & $\sim 2t^{3/2}/\sqrt{m}$   & V & Open \\
    4  & $2d^2/m$                  & $2t^2/m$                    & V & Open \\
    5  & $d^2/m + m$               & $t^2/m + m$                 & V & Open \\
    6  & $\sqrt{m^2+3d^2}$         & $\sqrt{3}\,t$               & V & Open \\
    7  & $d^2/m + d$               & $t^2/m + t$                 & V & Open \\
    8  & $\sqrt{d(m+3d)}$          & $\sqrt{3}\,t$               & V & Open \\
    9  & $(m+3d)/2$                & $3t/2$                      & V & Open \\
    10 & $\sqrt{d(d+3m)}$          & $t$                         & V & Open \\
    11 & $2m^3/d^2$                & $2m^3/t^2 \to 0$            & \textbf{I} & \textbf{Disproved} \\
    12 & $\sqrt{2m^2+2d^2}$        & $\sqrt{2}\,t$               & V & Open \\
    13 & $2m^4/d^3$                & $2m^4/t^3 \to 0$            & \textbf{I} & \textbf{Disproved} \\
    14 & $2d^3/m^2$                & $2t^3/m^2$                  & V & Open \\
    16 & $2d^4/m^3$                & $2t^4/m^3$                  & V & Open \\
    18 & $\sqrt{2m^3/d+2d^2}$      & $\sqrt{2}\,t$               & V & Open \\
    19 & $(4d^4{+}12dm^3)^{1/4}$   & $\sqrt{2}\,t$               & V & Open \\
    20 & $\tfrac{1}{2}\sqrt{7d^2{+}9m^2}$ & $\sqrt{7}\,t/2$      & V & Open \\
    21 & $\sqrt{d^3/m+3m^2}$       & $t^{3/2}/\sqrt{m}$          & V & Open \\
    22 & $(2d^4{+}14d^2m^2)^{1/4}$ & $\sqrt[4]{2}\,t$            & V & Open \\
    23 & $\sqrt{d^2+3dm}$          & $t$                         & V & Open \\
    24 & $(6d^4{+}10m^4)^{1/4}$    & $\sqrt[4]{6}\,t$            & V & Open \\
    25 & $(3d^4{+}13d^2m^2)^{1/4}$ & $\sqrt[4]{3}\,t$            & V & Open \\
    26 & $\tfrac{1}{2}\sqrt{5d^2{+}11dm}$ & $\sqrt{5}\,t/2$      & V & Open \\
    27 & $\sqrt{(3d^2{+}5dm)/2}$   & $\sqrt{3/2}\,t$             & V & Open \\
    30 & $m^3/d^2 + d^2/m$         & $t^2/m$                     & V & Open \\
    \bottomrule
  \end{tabular}
\end{table}

The hub-invisible mechanism explains why Bounds~11 and~13 are defeated:
as $t$ grows, the hub vertex---which dominates the Laplacian spectrum
via Proposition~\ref{prop:mu-soc}---becomes invisible to the bound formula.
The bound is then determined by the connector contribution, which
at $t = 2m - 1$ evaluates to~$2m$, while $\Lspec > 2m$ by
Proposition~\ref{prop:mu-soc}.
Conversely, the 22~hub-visible bounds grow at least as fast as the
spectral radius in~$t$, which is why our star-of-cliques families
cannot violate them.

\begin{remark}\label{rem:hub-invisible-unique}
  Bounds~11 and~13 are the \emph{only} vertex-type BHS bounds of the
  form $f(d,m) = 2m^p/d^{p-1}$ (with $p = 3, 4$).  All other vertex-type
  bounds contain terms that grow with~$d$, ensuring hub visibility.  This
  explains the pattern observed in Section~\ref{sec:main}: the
  star-of-cliques family violates precisely these two bounds and no others.
\end{remark}

\subsection{Edge-type mechanisms}\label{ssec:edge-mechanism}

The three edge-type bounds we disprove (40, 45, 48) fail by mechanisms
distinct from hub invisibility.

\subsubsection{Bound~40: transient violation window}\label{sssec:bound40}

Bound~40 is an edge-type bound violated by $\SoS(t, s)$.  Unlike the
divergent violations of Bounds~11 and~13, Bound~40 exhibits a
\emph{transient} violation: the gap $\Lspec - B_{40}$ is positive
only for $t$ in a bounded interval.

\begin{proposition}[Transient violation of Bound~40]\label{prop:bound40-transient}
  On $\SoS(t, s)$ with $s$ fixed, the violation gap of Bound~40 at the
  hub--intermediate edge satisfies:
  \begin{enumerate}[label=\textup{(\alph*)}]
    \item For $t \ll s^2$: $B_{40} > \Lspec$ \textup{(bound holds)}.
    \item For $t \sim s^2$: $B_{40} < \Lspec$ \textup{(bound fails)}.
    \item For $t \gg s^2$: $B_{40} > \Lspec$ \textup{(bound holds again)}.
  \end{enumerate}
  In particular, with $s = 8$ the maximum violation gap is approximately
  $1.27$ at $t = 39$.
\end{proposition}

\begin{proof}
  On $\SoS(t,s)$, the hub has degree $d_h = t$ and average neighbor degree
  $m_h = s + 1$ (each of its $t$ neighbors has degree $s+1$).  Each
  intermediate vertex has degree $d_w = s + 1$ and average neighbor degree
  $m_w = (t + s)/(s+1)$.

  The spectral radius satisfies
  $\Lspec(\SoS(t,s)) = \bigl(t + s + 2 + \sqrt{(t-s)^2 + 4s}\bigr)/2$
  (Proposition~\ref{prop:mu-sos}), which grows as $t + O(1)$ for large~$t$.

  The Bound~40 value at the hub--intermediate edge is
  \[
    B_{40} = 2 + \sqrt{2(m_h{-}1)^2 + 2(m_w{-}1)^2 + d_h^2 + d_w^2 - d_h m_h - d_w m_w}.
  \]
  Substituting the parameters and expanding, the dominant term under the
  square root for large~$t$ is $d_h^2 - d_h m_h = t^2 - t(s{+}1) = t(t{-}s{-}1)$,
  giving $B_{40} \sim 2 + t$ for $t \to \infty$.  Since
  $\Lspec \sim t + 1$, we have $B_{40} - \Lspec \sim 1 > 0$: the bound
  recovers.

  For small~$t$, both $\Lspec$ and $B_{40}$ are $O(s)$, and the bound
  also holds.  The violation occurs in a transient window near
  $t \approx s^2$ where the average neighbor degree term $m_w = (t{+}s)/(s{+}1)$
  is large enough to inflate $\Lspec$ via the quotient matrix mechanism
  but not yet large enough to dominate the bound formula.  The maximum
  violation gap of approximately~$1.27$ at $(t,s) = (39,8)$ was confirmed
  numerically.
\end{proof}

\subsubsection{Bound~45: antisymmetric mode decoupling}\label{sssec:bound45}

\begin{proposition}[Antisymmetric mode mechanism]\label{prop:bound45-mechanism}
  Bound~45 fails on $\DS(k,k)$ because the Laplacian spectral radius arises
  from the antisymmetric eigenvector sector, where the two hub vertices carry
  opposite signs.  The bound formula, which evaluates a symmetric expression
  at each edge, cannot detect this antisymmetric oscillation.
\end{proposition}

\begin{proof}
  By Proposition~\ref{prop:mu-ds}, the spectral radius of $\DS(k,k)$ is
  \[
    \Lspec\bigl(\DS(k,k)\bigr) = \frac{k + 3 + \sqrt{k^2 + 6k + 1}}{2},
  \]
  arising from the antisymmetric ($\sigma = -1$) sector of the equitable
  partition quotient.  The corresponding eigenvector has $x_{u} = -x_{v}$
  at the two hubs, and $x_{\ell_u} = -x_{\ell_v}$ at their respective
  pendant leaves.

  Bound~45 evaluates $g_{45}(d_i, d_j, m_i, m_j)$ at each edge.  On the
  hub--hub edge of $\DS(k,k)$, we have $d_i = d_j = k+1$ and
  $m_i = m_j = (k + k{+}1)/(k{+}1)$.  By symmetry, both edges
  incident to each hub yield the same value of~$g_{45}$,
  yet the actual spectral radius depends on the global
  topology---specifically the antisymmetric eigenmode across the
  hub--hub bridge.  The bound formula
  depends only on local degree parameters and cannot distinguish the
  symmetric eigenvalue
  $\lambda_{\mathrm{sym}} = (k{+}3 - \sqrt{k^2{+}6k{+}1})/2$ from the
  antisymmetric eigenvalue
  $\lambda_{\mathrm{anti}} = (k{+}3{+}\sqrt{k^2{+}6k{+}1})/2$.

  The violation gap satisfies
  $\Lspec - B_{45} = \Theta(1/k)$ as $k \to \infty$,
  so the gap vanishes asymptotically.  This suggests the bound is
  \emph{repairable} with an additive $O(1/\Delta)$ correction term.
\end{proof}

\subsubsection{Bound~48: degenerate small-graph failure}\label{sssec:bound48}

Bound~48 is violated only by~$P_3$.  Unlike the other four violations,
this appears to be a degenerate case: the bound formula contains a
denominator $2 + \sqrt{\cdots}$ that assumes nontrivial local structure.
On~$P_3$, with $d_i = 1$, $d_j = 2$ at the pendant--center edge, the
discriminant collapses and the bound evaluates to
$10/(2+\sqrt{2}) < 3 = \Lspec(P_3)$.  No counterexample was found
for $n \geq 4$ among all graphs tested (Section~\ref{ssec:exhaustive}),
and we conjecture that the bound holds for all connected graphs
with $n \geq 4$.

\subsection{Unified mechanism summary}\label{ssec:unified}

Table~\ref{tab:mechanism-summary} summarizes the failure mechanism
for each disproved bound.

\begin{table}[t]
  \centering
  \caption{Failure mechanisms for the five disproved BHS bounds.
    ``Gap behavior'' describes $\Lspec(G) - B(G)$ as graph parameters
    grow.  ``Repairable'' indicates whether an additive correction could
    restore validity.}
  \label{tab:mechanism-summary}
  {\small
  \begin{tabular}{@{}cllcl@{}}
    \toprule
    Bound & Type & Mechanism & Repairable? & Gap behavior \\
    \midrule
    11 & Vertex & Hub-invisible & No & $\to +\infty$ \\
    13 & Vertex & Hub-invisible & No & $\to +\infty$ \\
    40 & Edge & Transient window & Unclear & Bounded ($\max \approx 1.27$) \\
    45 & Edge & Antisym.\ decoupling & Yes & $\to 0$ as $\Theta(1/k)$ \\
    48 & Edge & Degenerate ($P_3$) & Yes & Fixed ($0.071$) \\
    \bottomrule
  \end{tabular}
  }
\end{table}

The two vertex-type failures (Bounds~11, 13) are \emph{irreparable}: the
violation gap diverges, so no bounded correction term can salvage these bounds.
Bound~45 is asymptotically repairable since the gap vanishes as
$O(1/\Delta)$.  Bound~48 is trivially repairable by excluding~$P_3$.
Bound~40 occupies an intermediate position: the violation is bounded but
occurs over a range of parameters, making a simple correction nontrivial.


%% =============================================================================
%% 6. CLASSIFICATION OF REMAINING BOUNDS
%% =============================================================================
\section{Status of the 38 open BHS bounds}\label{sec:classification}

Of the 68~bounds proposed by BHS~\cite{BHS2006}, Ghebleh et
al.~\cite{Ghebleh2024} disproved~30, leaving 38~open.
We have disproved~5 of these~38 (Section~\ref{sec:main}),
leaving \textbf{33~bounds open}.

\subsection{Classification criteria}\label{ssec:criteria}

We classify the 33~remaining open bounds into four categories based on
their behavior across all counterexample families tested ($\SoC(K_m, t)$
for $m \in [3,12]$, $t \in [2, 200]$; $\SoS(t, s)$ for $t \in [2, 100]$,
$s \in [1, 20]$; $\DS(k, k)$ for $k \in [1, 50]$; $\Cat(\ell, k)$ for
$\ell \in [2, 10]$, $k \in [0, 50]$; and all connected graphs with
$n \leq 10$):

\begin{definition}[Classification criteria]\label{def:classification}
  Let $\mathrm{gap}(B, G) = B(G) - \Lspec(G)$ denote the margin by
  which bound~$B$ holds on graph~$G$ (positive means the bound holds).
  \begin{enumerate}[label=\textup{(\roman*)}]
    \item \textbf{Super-safe}: The relative gap $\mathrm{gap}(B, G)/\Lspec(G)$
      exceeds~$0.1$ on all tested families (large relative margin).
    \item \textbf{Tight}: The infimum of $\mathrm{gap}(B, G)$ over all
      tested families equals zero (exact equality achieved on some graph
      family).
    \item \textbf{Safe}: The infimum of $\mathrm{gap}(B, G)/\Lspec(G)$
      over all tested families is bounded below by a positive constant
      (relative gap bounded away from zero).
    \item \textbf{Marginal}: Neither tight nor safe---the absolute gap is
      positive on all tested graphs but the relative gap approaches zero
      along some family.
  \end{enumerate}
\end{definition}

\subsection{Complete status table}\label{ssec:status-table}

Table~\ref{tab:classification} presents the complete status of all
38~bounds that were open after~\cite{Ghebleh2024}, organized by
resolution status.

\begin{table}[ht]
  \centering
  \caption{Status of the 38~BHS bounds left open by Ghebleh et
    al.~\cite{Ghebleh2024}.
    Bounds are grouped by resolution status.
    The 33~bounds that remain open are classified as
    \textsc{Super-safe}, \textsc{Tight}, \textsc{Safe}, or
    \textsc{Marginal} per Definition~\ref{def:classification}.}
  \label{tab:classification}
  {\small\renewcommand{\arraystretch}{0.88}%
  \begin{tabular}{lllll}
    \toprule
    Bound & Type & Resolved by & Method & Category \\
    \midrule
    \multicolumn{5}{l}{\emph{Disproved (5 bounds) --- this paper}} \\
    11  & Vertex & This paper & $\SoC(K_m, t)$    & Hub-invisible \\
    13  & Vertex & This paper & $\SoC(K_m, t)$    & Hub-invisible \\
    40  & Edge   & This paper & $\SoS(t, s)$      & Transient window \\
    45  & Edge   & This paper & $\DS(k, k)$       & Antisymmetric mode \\
    48  & Edge   & This paper & $P_3$             & Degenerate \\
    \midrule
    \multicolumn{5}{l}{\emph{Still open (33 bounds)}} \\
    1   & Vertex & ---         & --- & Super-safe \\
    4   & Vertex & ---         & --- & Super-safe \\
    5   & Vertex & ---         & --- & Super-safe \\
    6   & Vertex & ---         & --- & Super-safe \\
    7   & Vertex & ---         & --- & Super-safe \\
    8   & Vertex & ---         & --- & Safe \\
    9   & Vertex & ---         & --- & Super-safe \\
    10  & Vertex & ---         & --- & Marginal \\
    12  & Vertex & ---         & --- & Super-safe \\
    14  & Vertex & ---         & --- & Super-safe \\
    16  & Vertex & ---         & --- & Super-safe \\
    18  & Vertex & ---         & --- & Safe \\
    19  & Vertex & ---         & --- & Safe \\
    20  & Vertex & ---         & --- & Safe \\
    21  & Vertex & ---         & --- & Safe \\
    22  & Vertex & ---         & --- & Safe \\
    23  & Vertex & ---         & --- & Marginal \\
    24  & Vertex & ---         & --- & Safe \\
    25  & Vertex & ---         & --- & Safe \\
    26  & Vertex & ---         & --- & Safe \\
    27  & Vertex & ---         & --- & Safe \\
    30  & Vertex & ---         & --- & Safe \\
    33  & Edge   & ---         & --- & Tight \\
    34  & Edge   & ---         & --- & Safe \\
    35  & Edge   & ---         & --- & Safe \\
    37  & Edge   & ---         & --- & Safe \\
    38  & Edge   & ---         & --- & Safe \\
    39  & Edge   & ---         & --- & Safe \\
    42  & Edge   & ---         & --- & Tight \\
    44  & Edge   & ---         & --- & Safe \\
    46  & Edge   & ---         & --- & Safe \\
    47  & Edge   & ---         & --- & Tight \\
    56  & Edge   & ---         & --- & Safe \\
    \bottomrule
  \end{tabular}%
  }
\end{table}

\noindent
Of the 33~remaining open bounds, all vertex-type ones are hub-visible
(Proposition~\ref{prop:hub-classification}).  Three edge-type bounds
(33, 42, and~47) are \emph{tight}: they achieve exact equality on complete
bipartite graphs and windmill graphs.

\begin{remark}[Comparison with Ghebleh et al.]\label{rem:ghebleh-comparison}
  Of the 30~bounds disproved by Ghebleh et al.~\cite{Ghebleh2024},
  all counterexamples have at most 13~vertices.  None of the 38~bounds
  left open by their search can be disproved by any graph with $n \leq 13$,
  and our smallest new counterexample ($P_3$ aside, which has $n = 3$) has
  $n = 14$ (for Bound~45).  The largest has $n = 352$ (for Bound~40).
  This suggests a phase transition in counterexample difficulty: the
  ``easy'' bounds are disproved by small graphs findable by stochastic
  search, while the ``hard'' ones require structured infinite families
  with closed-form spectral analysis.
\end{remark}

\subsection{Exhaustive verification of the tight bounds}\label{ssec:exhaustive}

We performed exhaustive computational verification of the three
tight bounds (33, 42, and~47) on every connected graph with $n \leq 10$
vertices---a total of $11{,}988{,}768$ graphs.
Zero violations were found.

\begin{proposition}[Finite verification]\label{thm:finite-verification}
  Bounds~33, 42, and~47 from~\cite{BHS2006} hold for all connected graphs
  with at most~$10$ vertices.
\end{proposition}

Windmill graphs $W(k,3)$ (the friendship graphs) provide asymptotically
tight instances for Bounds~42 and~47: the relative gap satisfies
$\bigl(B_{42} - \Lspec\bigr)/\Lspec \sim \sqrt{3}/k \to 0$ as
$k \to \infty$.  Proving or disproving these three bounds remains the
most compelling open problem arising from this work.


%% =============================================================================
%% 7. DISCUSSION
%% =============================================================================
\section{Discussion}\label{sec:discussion}

\paragraph{Summary.}
BHS~\cite{BHS2006} proposed 68~conjectured bounds for~$\Lspec(G)$.
Ghebleh et al.~\cite{Ghebleh2024} disproved~30 via cross-entropy methods.
We have disproved~5 more by structural construction.
Together, 35 of the original 68~bounds are now resolved (52\%),
with 33~remaining open.

\paragraph{Structural vs.\ computational approaches.}
Cross-entropy search~\cite{Ghebleh2024,Wagner2021,BouffardBreen2025}
explores adjacency matrices for $n \leq 24$.
Our structural approach targets algebraic weaknesses
(hub-invisible ratios, antisymmetric modes)
and produces infinite families with
closed-form spectral radii.  The two approaches are complementary:
stochastic search excels at finding small counterexamples, while
structural construction reaches the large-$n$ regime ($n = 352$ for
Bound~40) where no counterexample can be found computationally.

\paragraph{Gap asymptotics and repairability.}
The five violations exhibit qualitatively different gap behaviors
(Table~\ref{tab:mechanism-summary}).
Bounds~11 and~13 have divergent gaps (irreparable).
Bound~40 has a bounded transient window (max gap $\approx 1.27$).
Bound~45 has a vanishing gap $\Theta(1/k)$, suggesting repairability
with an additive $O(1/\Delta)$ correction.
Bound~48 fails only on~$P_3$.
Understanding which bounds are repairable may guide the design of
improved automated conjectures.

\paragraph{A weighted Laplacian intermediate result.}
In attempting to prove the tight bounds (42, 47), we obtained the
following intermediate result, which we include here as a potentially
useful tool for future work.

\begin{proposition}[Weighted Laplacian bound]\label{prop:weighted-laplacian}
  For any connected graph~$G$,
  \[
    \Lspec(G)^2 \leq \lambda_{\max}(M),
  \]
  where $M$ is the Laplacian of the weighted graph on the same vertex set
  with edge weights $w_{ij} = d_i + d_j$ for each $\{i,j\} \in E$.
\end{proposition}

\begin{proof}[Proof sketch]
  Let $x$ be the eigenvector of~$L(G)$ with eigenvalue~$\Lspec(G)$, normalized
  so that $\|x\|_2 = 1$. The eigenequation gives
  $\Lspec(G)\, x_v = d_v x_v - \sum_{u \sim v} x_u$. Squaring both sides:
  \[
    \Lspec(G)^2\, x_v^2 = \Bigl(d_v x_v - \sum_{u \sim v} x_u\Bigr)^{\!2}.
  \]
  Expanding the square gives $d_v^2 x_v^2 - 2 d_v x_v \sum_{u \sim v} x_u
  + \bigl(\sum_{u \sim v} x_u\bigr)^2$.
  Applying Cauchy--Schwarz to the cross-term,
  $-2 d_v x_v \sum_{u \sim v} x_u \leq d_v \sum_{u \sim v}(x_v^2 + x_u^2)$,
  and to $\bigl(\sum_{u \sim v} x_u\bigr)^2 \leq
  d_v \sum_{u \sim v} x_u^2$, then summing over all
  vertices yields $\Lspec(G)^2 \leq x^T M\, x$, where $M$ has entries
  \[
    M_{vv} = \sum_{u \sim v} (d_v + d_u), \qquad
    M_{vu} = -(d_v + d_u) \;\text{ if } u \sim v.
  \]
  The Rayleigh quotient bound gives $x^T M\, x \leq \lambda_{\max}(M)$.
\end{proof}

\begin{remark}\label{rem:weighted-laplacian-gap}
  Proposition~\ref{prop:weighted-laplacian} implies that proving Bound~42,
  for example, reduces to showing
  $\lambda_{\max}(M) \leq \max_{(i,j) \in E}(d_i^2 + d_j^2 + 2m_i \cdot m_j)$.
  This second inequality remains open.
\end{remark}

\paragraph{Proof attempts for tight bounds.}
We attempted to prove the two remaining tight bounds (42, 47) using Rayleigh
quotient methods, Schur complement analysis, Gershgorin circles, and degree
sequence constraints.  While the weighted Laplacian bound
(Proposition~\ref{prop:weighted-laplacian}) provides a useful intermediate
step, complete proofs remain elusive.
The fundamental obstacle is that standard spectral techniques do not
naturally produce edge-maximum bounds involving average neighbor degrees:
the Rayleigh quotient $x^T L x / x^T x$ yields vertex-sum expressions,
whereas edge-type bounds require converting these into edge-maximum form.

\paragraph{Open questions.}
\begin{enumerate}[label=(\alph*)]
  \item Can the three tight bounds (33, 42, 47) be proved?
  \item Can additional bounds among the remaining~33 be proved?
  \item Do three-level hierarchies (e.g., star-of-stars-of-cliques)
    violate additional bounds?
  \item Can Bound~45 be repaired with an $O(1/\Delta)$ correction?
  \item What is the correct minimum-order threshold for Bound~48?
\end{enumerate}


%% =============================================================================
%% AI DISCLOSURE (Required by Elsevier)
%% =============================================================================
\section*{Declaration of Generative AI in Scientific Writing}

During the preparation of this work, the author used Claude
(Anthropic; model: Claude Opus 4.6) for computational assistance,
including code development for graph construction and eigenvalue
computation, structural analysis of graph families, and systematic
counterexample search. The author reviewed and edited all content and
takes full responsibility for the publication.


%% =============================================================================
%% ACKNOWLEDGMENTS
%% =============================================================================
\section*{Acknowledgments}

The author thanks Pierre Hansen and Dragan Stevanovi\'{c} for making the
BHS bound list publicly available.  Computational assistance was provided
by Claude (Anthropic); all mathematical results were independently verified
by the author.
%% Acknowledge Claude's computational role.


%% =============================================================================
%% REFERENCES
%% =============================================================================
\bibliographystyle{elsarticle-num}
\bibliography{references}

\end{document}
